% =====
\documentclass[11pt,a4paper]{article}

% Seitenraender (2 cm)
\usepackage[a4paper,top=2cm,bottom=2cm,left=2cm,right=2cm]{geometry}

% Kodierung + Sprache: Deutsch (pdfLaTeX)
\usepackage[T1]{fontenc}
\usepackage[utf8]{inputenc}
\usepackage[ngerman]{babel}

% Schrift: Arial-Ersatz (Helvetica-Klon), pdfLaTeX-tauglich und schnell
\usepackage{helvet}
\renewcommand{\familydefault}{\sfdefault}

% Zeilenabstand 1,5
\usepackage{setspace}
\onehalfspacing

% Blocksatz + Mikrotypografie
\usepackage{microtype}
\setlength{\parindent}{0pt}
\setlength{\parskip}{6pt}

% Ueberschriften ~12pt fett
\usepackage{titlesec}
\titleformat{\section}{\fontsize{12}{14}\bfseries}{\thesection}{1em}{}
\titleformat{\subsection}{\fontsize{12}{14}\bfseries}{\thesubsection}{1em}{}
\titleformat{\subsubsection}{\fontsize{12}{14}\bfseries}{\thesubsubsection}{1em}{}

% Max. 3 Kapitelstufen
\setcounter{secnumdepth}{3}
\setcounter{tocdepth}{3}

% Fussnoten 10pt
\usepackage{footmisc}
\renewcommand{\footnotesize}{\fontsize{10}{12}\selectfont}

% Zitation: APA (author-year) via natbib + apalike
\usepackage[round,authoryear]{natbib}
\setcitestyle{authoryear,round}

% Seitenzahlen unten zentriert
\usepackage{fancyhdr}
\pagestyle{fancy}
\fancyhf{}
\cfoot{\thepage}
\renewcommand{\headrulewidth}{0pt}

% Grafiken + Tabellen + Links
\usepackage{graphicx}
\graphicspath{{./}{figures/}}
\usepackage{float}
\usepackage{booktabs}
\usepackage{array}
\usepackage{ragged2e}
\usepackage{ltablex} % tabularx, das ueber Seiten umbrechen kann
\keepXColumns
\usepackage{amssymb}
\usepackage{hyperref}
\setcounter{LTchunksize}{200}

% Umbrechbare, links ausgerichtete X-Spalte
\newcolumntype{Y}{>{\RaggedRight\arraybackslash}X}

\begin{document}

% ---------------- Titelblatt (keine Seitenzahl) ----------------
\begin{titlepage}
\thispagestyle{empty}
\centering
\vspace*{2.5cm}

{\LARGE\bfseries KONFIGURATION DER SOFTWAREENTWICKLUNG\\[0.2cm] UND QUALITÄTSPLANUNG\par}
\vspace{0.6cm}
{\Large Meilenstein 3\par}
\vspace{1.2cm}

{\Large ISEF01: Projekt Software Engineering\par}
\vspace{1.5cm}

{\large Thema E\par}
\vspace{0.2cm}
{\Large\bfseries Geschenke-Manager\par}
\vspace{0.3cm}
{\large Eine datenschutzorientierte, browserbasierte Webanwendung\par}

\vfill

{\large
\begin{tabular}{@{}ll@{}}
Julian Ammann & 31911823 \\
Yin Yin Wu-Hanke & 42303256 \\
Anton Hirsch & IU14119299 \\
Kevin Jordan Taghu & IU14150888 \\
\end{tabular}\par}

\vspace{1.0cm}
{\large
Tutor: Prof.\ Dr.\ rer.\ nat.\ Tobias Brückmann\par
Datum: 28.\,September 2026
\par}

\end{titlepage}

% ---------------- Vorderteil: roemisch ----------------
\pagenumbering{Roman}
\setcounter{page}{1}

\tableofcontents

% ---------------- Textteil: arabisch ----------------
\clearpage
\pagenumbering{arabic}

\section{Konfiguration der Softwareentwicklung}

\subsection{Ausgangslage und Ziel}

Dieses Kapitel legt fest, wie das Team den Geschenke-Manager bis MS 4 entwickelt, prüft und bereitstellt. Es beschreibt den geplanten Zielzustand. Alle Festlegungen gelten ab MS 3 verbindlich für alle Teammitglieder.

Zu MS 3 besteht ein technisches Grundgerüst: ein gemeinsames Repository für Backend und Frontend, gepinnte Werkzeugversionen, Git-Hooks, CI-Workflows, Dockerfiles und Health-Endpunkte. Die übrigen technischen und fachlichen Funktionen sind geplant und entstehen bis MS 4.

\subsection{Vorgehensmodell}

Das Team arbeitet iterativ in vier Abschnitten. Jeder Abschnitt endet mit einem lauffähigen Stand auf \texttt{main}.

\begin{tabularx}{\linewidth}{@{}l >{\hsize=0.60\hsize\RaggedRight\arraybackslash}X >{\hsize=1.40\hsize\RaggedRight\arraybackslash}X@{}}
\toprule
\textbf{Abschnitt} & \textbf{Zeitraum} & \textbf{Ziel} \\
\midrule
\endhead
Sprint 0 & bis 27.09.2026 (MS 3) & Entwicklungskonfiguration, verbindliches Datenmodell, Authentifizierung, leere Anwendung (Walking Skeleton) unter HTTPS erreichbar \\
Sprint 1 & 28.09. bis 03.10.2026 & Mandantentrennung als wiederverwendbares Muster, erste Fachdomäne „Personen und Anlässe“ durch alle Schichten bis ins Frontend \\
Sprint 2 & 04.10. bis 08.10.2026 & Geschenke, Beschenkungen, Aufgaben, Notizen, Anhänge, Benachrichtigungen, Teilen-Links, HTML-Ansicht und Geschenkvorschläge parallel \\
Stabilisierung & 09.10. bis 11.10.2026 & Feature-Freeze ab 08.10.2026. Nur noch Fehlerbehebung, Tests, Dokumentation und Bereitstellung für MS 4 \\
\bottomrule
\end{tabularx}

Die erste Fachdomäne dient als Vorlage für alle weiteren. Erst wenn sie mit Tests und Frontend-Anbindung fertig ist, beginnen die übrigen Domänen parallel, und alle vier Entwickler schreiben ihre Modelle und Migrationen nach diesem erprobten Muster.

Jeder Abschnitt beginnt mit einem Regelmeeting in Microsoft Teams. Dort verteilt das Team die Aufgaben und prüft den Fortschritt gegen den Projektstrukturplan. Kurzfristige Blocker klärt das Team über Signal. Abnahmekriterien jedes Abschnitts sind die Qualitätsziele QZ-01 bis QZ-09 und die Definition of Done am Ende dieses Kapitels.

\subsection{Anforderungsmanagement}

Grundlage sind die Anforderungen aus den in MS 1 festgelegten funktionalen Anforderungen F-01 bis F-19 und Qualitätsanforderungen Q-01 bis Q-09.

\textbf{Priorisierung}: Die Priorisierung aus MS 1 gilt weiterhin. MUSS kennzeichnet dabei den verbindlichen Kernumfang (F-01 bis F-17). SOLL kennzeichnet Erweiterungen (F-18, F-19), die nur umgesetzt werden, wenn der Kernumfang und die Qualität dadurch nicht gefährdet sind.

\textbf{Konkretisierung}: Die Anforderungen aus MS 1 sind bewusst auf grober Ebene formuliert. Das Team konkretisiert sie schrittweise während der Umsetzung, jeweils zu Beginn des Abschnitts, in dem sie umgesetzt werden. Dabei werden sie in umsetzbare Aufgaben zerlegt und um prüfbare Akzeptanzkriterien ergänzt. Aus den Akzeptanzkriterien entstehen die Testfälle. Offene fachliche Fragen klärt das Team in der Planungphase. Fragen, die nur der Tutor beantworten kann, stellt die Projektleitung über Redmine.

\textbf{Nachvollziehbarkeit}: Die Kennungen aus MS 1 bleiben über das gesamte Projekt hinweg erhalten. Aufgaben, Änderungen und Tests verweisen auf die betroffenen Anforderungen. Aufgaben werden als GitHub Issues im Projekt-Repository geführt und im GitHub Project „Geschenke-Manager Backlog“ geplant. Jedes Issue nennt im Titel oder im Feld „Anforderung“ die Anforderungs-ID und enthält die Akzeptanzkriterien als Checkliste; die PSP-Nummer steht in der Beschreibung. Das Project führt je Issue die Felder Abschnitt, Zieltermin, Priorität (MUSS oder SOLL), Bereich und Status. Die verantwortliche Person wird dem Issue zugewiesen. Das Product Backlog umfasst alle offenen Issues des Projects; das Sprint Backlog sind die Issues, deren Feld „Abschnitt“ auf den laufenden Abschnitt gesetzt ist. Die Bugliste sind die Issues mit dem Label \texttt{bug}; sie enthalten Reproduktionsschritte sowie Verweise auf Anforderung und Pull Request. Blocker werden ebenfalls als Issue erfasst und zusätzlich über Signal gemeldet. Pull Requests schließen ihre Issues über \texttt{Closes \#<Nummer>}. Testfälle liegen in \texttt{backend/\allowbreak{}tests/\allowbreak{}} beziehungsweise den Frontend-Tests und führen die geprüften Anforderungs-IDs. Die Rückverfolgbarkeitsmatrix aus dem Testlauf (QZ-01) zeigt, ob eine Anforderung erfüllt ist. Testprotokolle werden aus den CI-Läufen übernommen und im Testabschlussbericht dokumentiert. Wesentliche Projektentscheidungen werden als GitHub Issue mit Begründung festgehalten, die Prüfungen der Liefergegenstände samt Ergebnis im Pull Request oder am zugehörigen Redmine-Ticket. Damit wird Q-09 abgedeckt.

\textbf{Qualitätssicherung}: Bevor eine Anforderung umgesetzt wird, wird die Konkretisierung darauf geprüft, ob sie eindeutig, testbar und widerspruchsfrei ist und ob Akzeptanzkriterien vorliegen. Nach der Umsetzung wird die fachliche Korrektheit gegen die Anforderung geprüft.

\textbf{Änderungen}: Neue oder geänderte Anforderungen bewertet das Team in der nächsten Planung und dokumentiert die Entscheidung nachvollziehbar. Reicht die Kapazität nicht aus, gilt die Reihenfolge aus dem Risikomanagement in MS 1 (PM-02, PM-03, PM-07):

\begin{enumerate}
\item SOLL-Anforderungen werden zurückgestellt.
\item Vom Team gewählte Ergänzungen werden vereinfacht oder als SOLL eingeordnet.
\item Die Anforderungen des Themenblatts bleiben verbindlich.
\end{enumerate}


Ab dem Feature-Freeze am 08.10.2026 werden keine neuen Anforderungen mehr aufgenommen.

\subsection{Rollen und Zuständigkeiten in der Entwicklung}

\begin{tabularx}{\linewidth}{@{}>{\hsize=0.47\hsize\RaggedRight\arraybackslash}X >{\hsize=1.24\hsize\RaggedRight\arraybackslash}X >{\hsize=1.30\hsize\RaggedRight\arraybackslash}X@{}}
\toprule
\textbf{Teammitglied} & \textbf{Entwicklung} & \textbf{Review und Prüfung} \\
\midrule
\endhead
Julian Ammann & Einrichtung von Repository, Werkzeugen, CI/CD, Containern und Hosting (PSP 4.1). Authentifizierung und Sitzungen (PSP 4.2). Frontend & Reviewt Backend-Pull-Requests sowie Frontend-Pull-Requests, die er nicht selbst erstellt hat. Pflicht-Reviewer für Container-, Workflow- und Deploy-Konfiguration. Führt die Deploys aus \\
Kevin Jordan Taghu & Backend: Personen und Anlässe (PSP 4.3), Benachrichtigungen mit Scheduler und E-Mail-Versand (PSP 4.5), Teilen-Links und HTML-Ansicht (PSP 4.6) & Reviewt die Backend-Pull-Requests von Anton Hirsch. Koordiniert als Projektleitung die Abschnitte \\
Anton Hirsch & Backend: Geschenke, Beschenkungen, Aufgaben, Notizen und Anhänge (PSP 4.4), Geschenkvorschläge (PSP 4.7). Fachliche Testfälle & Reviewt die Backend-Pull-Requests von Kevin Jordan Taghu. Prüft die fachliche Korrektheit gegen die Anforderungen \\
Yin Yin Wu-Hanke & Frontend, Berechtigungs- und Sicherheitstests, Auswertung der Secret- und Abhängigkeitsscans & Reviewt Frontend-Pull-Requests, die sie nicht selbst erstellt hat. Prüft sicherheitsrelevante Änderungen (Anmeldung, Teilen-Links, Uploads) \\
\bottomrule
\end{tabularx}

Jeder Pull Request braucht die Freigabe einer anderen Person. Backend-Code prüft Julian Ammann oder der jeweils andere Backend-Entwickler. Frontend-Code prüft Julian Ammann oder Yin Yin Wu-Hanke, jeweils nicht die erstellende Person. Sicherheitsrelevante Änderungen an Anmeldung, Teilen-Links und Uploads prüft Yin Yin Wu-Hanke zusätzlich aus Sicherheitssicht.

Die Backend-Stränge sind nach Fachdomänen getrennt. Jede Domäne erhält in jeder Schicht eine eigene Datei, etwa für Personen je ein Modul für Modell, Schema, Service und Router. Gemeinsam bearbeitet werden nur die beiden Dateien, in denen Modelle und Router registriert werden. Nur dort können Merge-Konflikte entstehen.

Backend und Frontend verbindet die OpenAPI-Spezifikation, die das Backend aus dem Code erzeugt. Ein Endpunkt gilt als übergeben, sobald er mit Zusammenfassung, Beschreibung und Beispielen in der Spezifikation steht. Das Frontend verwendet ausschließlich die daraus generierten Typen.

\subsection{Architektur}

{\footnotesize
\begin{verbatim}
Browser --HTTPS--> Traefik --> Frontend (SvelteKit, Node, :3000)
                                    |
                                    | serverseitige Datenabfragen und Formulare
                                    | Proxy für Browser-Aufrufe auf /api/*
                                    v
                                Backend (FastAPI, :8000) --> PostgreSQL 17
                                    ^                          ^
                                    |                          |
                                Scheduler (gleiches Image) ----+
                                    |
                                    v
                                SMTP-Dienst
\end{verbatim}
}

\begin{itemize}
\item Öffentlich erreichbar ist nur das Frontend. Backend, Scheduler und Datenbank liegen in einem internen Container-Netz.
\item Der Browser sieht nur eine Adresse. Das Session-Cookie ist deshalb ein First-Party-Cookie, und eine CORS-Freigabe entfällt.
\item Der Scheduler läuft als eigener Prozess aus dem Backend-Image, immer mit genau einer Instanz. Im API-Prozess läuft kein Scheduler.
\item Bilddateien liegen auf einem eigenen Volume und werden über eine geschützte API-Route ausgeliefert: mit gültiger Sitzung oder, soweit ausdrücklich freigegeben, mit gültigem Freigabe-Link.
\end{itemize}


\subsection{Technologie-Stack}

\begin{tabularx}{\linewidth}{@{}>{\hsize=0.91\hsize\RaggedRight\arraybackslash}X >{\hsize=1.37\hsize\RaggedRight\arraybackslash}X >{\hsize=0.72\hsize\RaggedRight\arraybackslash}X@{}}
\toprule
\textbf{Bereich} & \textbf{Festlegung} & \textbf{Version} \\
\midrule
\endhead
Backend-Sprache & Python & 3.14 \\
Web-Framework & FastAPI mit Uvicorn, vollständig asynchron & $\geq$ 0.141 \\
Datenzugriff & SQLAlchemy 2.0 (async) mit asyncpg & 2.x \\
Datenbank & PostgreSQL & 17 \\
Schemamigrationen & Alembic & $\geq$ 1.20 \\
Validierung und Konfiguration & Pydantic, pydantic-settings & 2.x \\
Passwort-Hashing & Argon2id & - \\
Logging & structlog, lesbar in der Entwicklung, JSON in Produktion & $\geq$ 26.1 \\
Scheduler & APScheduler in eigenem Prozess & - \\
E-Mail & aiosmtplib, lokal Mailpit als Test-Postfach & - \\
Bildverarbeitung & Pillow für Formatprüfung und Entfernung von EXIF-Daten & - \\
Frontend-Sprache & TypeScript & 6.x \\
Frontend-Framework & SvelteKit 2 mit Svelte 5, serverseitiges Rendering auf Node & 2.x / 5.x \\
Styling & Tailwind CSS & 4.x \\
API-Client & \texttt{openapi-fetch} mit aus der Spezifikation generierten Typen & 0.17 / 7.x \\
Laufzeit Frontend & Node.js & 26 \\
Paketverwaltung & uv (Backend), pnpm (Frontend) & 0.12.9 / 12.5.1 \\
Tests & pytest, httpx, testcontainers, Vitest, Playwright & - \\
\bottomrule
\end{tabularx}

Python 3.14, Node 26 und TypeScript 6 sind sehr neue Versionen. Lockfiles und gepinnte Werkzeuge begrenzen das Risiko. Tritt ein blockierender Fehler auf, wechselt das Team für das betroffene Werkzeug auf die vorige Major-Version.

\subsection{Repository-Struktur}

{\footnotesize
\begin{verbatim}
isefgm/
+-- backend
|   +-- alembic/versions      Migrationen
|   +-- src/giftmanager
|   |   +-- api/v1            Router und Endpunkte
|   |   +-- core              Konfiguration, Datenbank, Anmeldung, Logging
|   |   +-- models            Datenbankmodelle
|   |   +-- schemas           Request- und Response-Schemas
|   |   +-- services          Geschäftslogik
|   +-- tests
+-- frontend
|   +-- src/lib/api           API-Client und generierte Typen
|   +-- src/routes            Seiten und Proxy für /api/*
+-- deploy                    Produktionskonfiguration und Betriebsanleitung
+-- docs                      Meilenstein-Dokumente, Datenmodell, generierte Dokumentation
\end{verbatim}
}

\subsection{Konfiguration des Backends}

\subsubsection{Schichten}

\begin{tabularx}{\linewidth}{@{}l >{\hsize=1.10\hsize\RaggedRight\arraybackslash}X >{\hsize=0.90\hsize\RaggedRight\arraybackslash}X@{}}
\toprule
\textbf{Schicht} & \textbf{Verantwortung} & \textbf{Darf nicht} \\
\midrule
\endhead
API & HTTP-Routing, Anmeldeprüfung, Request- und Response-Schemas, Statuscodes & Geschäftslogik enthalten oder SQL formulieren \\
Schemas & Validierung und Serialisierung, mit Beschreibungen und Beispielen & auf die Datenbank zugreifen \\
Services & Geschäftsregeln, Besitzerprüfung, Datenbankabfragen & HTTP-Fehler werfen, Transaktionen abschließen, die Uhrzeit abfragen \\
Modelle & Tabellen, Beziehungen, Constraints, Kommentare & Logik enthalten \\
\bottomrule
\end{tabularx}

\subsubsection{Verbindliche Regeln}

Diese Regeln prüft das Team in jedem Backend-Review.

\begin{enumerate}
\item \textbf{Transaktion.} Jeder Request läuft in genau einer Datenbanktransaktion. Endet der Endpunkt normal, wird sie bestätigt, bei einem Fehler zurückgerollt. Services schließen keine Transaktionen selbst ab. Der Scheduler verwendet eine eigene Transaktion je Nachricht.
\item \textbf{Mandantentrennung (QZ-06).}
\begin{itemize}
\item Jede Tabelle mit Nutzerdaten erhält eine Spalte \texttt{owner\_\allowbreak{}id} mit Fremdschlüssel auf das Konto. Wird das Konto gelöscht, werden die Daten mitgelöscht.
\item Services erhalten \texttt{owner\_\allowbreak{}id} als ersten Parameter und filtern jede Abfrage danach. Ein Zugriff nur über die ID ohne Besitzerfilter ist verboten.
\item Fremde IDs beantwortet die API mit \texttt{404}, nie mit \texttt{403}, weil ein \texttt{403} verraten würde, dass der Datensatz existiert.
\item Beim Verknüpfen zweier Datensätze prüft der Service, dass beide demselben Konto gehören. Ausgenommen sind die systemweiten Anlasstypen.
\item Zu jeder Ressource gehört ein Test „Nutzer B greift auf Daten von Nutzer A zu $\rightarrow$ 404“.
\end{itemize}

\item \textbf{Fehler.} Fachliche Fehler haben eine gemeinsame Basisklasse. Ein zentraler Handler übersetzt sie, Validierungsfehler und HTTP-Fehler in ein einheitliches Format nach RFC 9457 (Problem Details) mit Typ, Titel, Status, Beschreibung und optional einer Liste von Feldfehlern.
\item \textbf{Listen (QZ-02).} Jede Liste ist paginiert, standardmäßig 50 und höchstens 200 Einträge pro Seite. Die Antwort enthält die Einträge, die Gesamtzahl und die Seitenparameter. Beziehungen werden gebündelt geladen, nicht einzeln in Schleifen.
\item \textbf{Zeit.} Zeitpunkte werden in UTC gespeichert, Kalenderdaten als reines Datum. Die aktuelle Zeit liefert eine austauschbare Uhr, in Tests etwa ein festes Datum für Weihnachtsbenachrichtigungen. Kalenderauswertungen laufen in der Zeitzone Europe/Berlin.
\item \textbf{Dokumentation im Code.} Jede Tabelle und Spalte erhält einen Datenbankkommentar. Jeder Endpunkt erhält Zusammenfassung und Beschreibung, jedes Schema-Feld Beschreibung und Beispiel. Jeder Test nennt die Anforderungs-IDs, die er prüft. Diese Texte sind Deutsch, Code und Code-Kommentare Englisch.
\end{enumerate}


\subsubsection{API-Konventionen}

\begin{itemize}
\item Alle Endpunkte liegen unter \texttt{/\allowbreak{}api/\allowbreak{}v1}. Ressourcen heißen im Plural, zum Beispiel \texttt{/\allowbreak{}persons}, \texttt{/\allowbreak{}gifts}, \texttt{/\allowbreak{}giftings}, \texttt{/\allowbreak{}tasks}.
\item Jeder Endpunkt hat eine eindeutige Operation-ID in camelCase, etwa \texttt{listPersons} oder \texttt{createGift}. Daraus entstehen die Methodennamen im Frontend-Client.
\item Schreibende Endpunkte geben die vollständige Ressource zurück. Nur Löschungen antworten mit \texttt{204} ohne Inhalt.
\item Statuswerte sind englische Schlüssel wie \texttt{idea}, \texttt{planned}, \texttt{acquired}, \texttt{given}. Die deutsche Anzeige übernimmt das Frontend.
\item Kalenderdaten haben das Format \texttt{YYYY-MM-DD}, Zeitpunkte ISO 8601 mit Zeitzone.
\end{itemize}


\subsubsection{Authentifizierung und Sitzungen}

Alle kontogebundenen fachlichen Endpunkte setzen eine Anmeldung voraus. Ausnahmen sind Registrierung, Anmeldung und die Nur-Lese-Ansicht über einen gültigen Freigabe-Link. Diese Ansichten geben ausschließlich die ausdrücklich freigegebenen Personen und Geschenkideen zurück; Schreibzugriffe bleiben ausgeschlossen. Julian Ammann setzt die Authentifizierung in Sprint 0 um.

\begin{tabularx}{\linewidth}{@{}>{\hsize=0.53\hsize\RaggedRight\arraybackslash}X >{\hsize=1.47\hsize\RaggedRight\arraybackslash}X@{}}
\toprule
\textbf{Aspekt} & \textbf{Festlegung} \\
\midrule
\endhead
Passwörter & gehasht mit Argon2id, Länge 8 bis 128 Zeichen \\
Anmeldename & E-Mail-Adresse, in Kleinbuchstaben gespeichert, eindeutig \\
Sitzung & eigene Tabelle mit Erstellungs-, Ablauf- und letztem Nutzungszeitpunkt. Die Datenbank speichert nur einen Hash des Sitzungstokens \\
Cookie & \texttt{HttpOnly}, \texttt{SameSite=Lax}, außerhalb der Entwicklung \texttt{Secure}. Laufzeit 14 Tage, Verlängerung bei Nutzung \\
Schutz vor Ausspähen & Auch bei unbekannter E-Mail wird ein Passwort-Hash geprüft. Die Antwortzeit verrät nicht, ob ein Konto existiert \\
Rate-Limit & höchstens 10 Anmelde- oder Registrierungsversuche je E-Mail-Adresse in 15 Minuten \\
CSRF-Schutz & Zusätzlich zu \texttt{SameSite} prüft das Backend bei schreibenden Aufrufen, ob die Anfrage von der eigenen Adresse stammt \\
Endpunkte & Registrieren, Anmelden, Abmelden, eigenes Konto abrufen, Konto löschen (F-17) \\
Kontolöschung & löscht das Konto samt Sitzungen, allen zugehörigen Daten und hochgeladenen Bildern \\
\bottomrule
\end{tabularx}

Automatisierte Tests decken mindestens ab: Registrierung, doppelte E-Mail, zu kurzes Passwort, falsches Passwort, unbekanntes Konto, Anmeldung, Abmeldung, abgelaufene Sitzung, Rate-Limit, Kontolöschung und Zugriff ohne Anmeldung.

\subsubsection{Datenmodell und Migrationen}

Das Team legt in Sprint 0 ein vollständig attributiertes Datenmodell als ER-Diagramm fest. Es baut auf dem Zielmodell aus MS 1 auf und ergänzt die zentrale Entität Beschenkung (\texttt{gifting}), die Zuordnungstabellen, Notizen, das Nutzerprofil und eine Kategorie für Geschenke als Grundlage der Vorschläge (F-16). Ab MS 3 ist dieses Modell für alle Entwickler verbindlich.

Konventionen:

\begin{itemize}
\item Primärschlüssel \texttt{id} als UUIDv7.
\item Die Kontotabelle heißt \texttt{user\_\allowbreak{}account}, weil \texttt{user} in PostgreSQL reserviert ist.
\item Statuswerte werden als Text gespeichert und per \texttt{CHECK}-Constraint auf erlaubte Werte begrenzt: Beschenkung \texttt{idea | planned | acquired | given}, Aufgabe \texttt{open | in\_\allowbreak{}progress | done | discarded}, Benachrichtigung \texttt{planned | sending | sent | failed}.
\item Constraint-Namen folgen einer festen Konvention (\texttt{pk\_\allowbreak{}}, \texttt{fk\_\allowbreak{}}, \texttt{uq\_\allowbreak{}}, \texttt{ix\_\allowbreak{}}, \texttt{ck\_\allowbreak{}}), damit Migrationen reproduzierbar bleiben.
\item Wiederkehrende Anlässe speichern die Regel, die Beschenkung speichert den konkreten Termin.
\item Die systemweiten Anlasstypen „Geburtstag“ und „Weihnachten“ legt eine Migration an.
\end{itemize}


Ablauf einer Schemaänderung:

\begin{enumerate}
\item Modell in der Datei der eigenen Domäne anlegen und registrieren.
\item Migration automatisch aus dem Modell erzeugen.
\item Migration von Hand gegenlesen, weil die automatische Erzeugung keine Umbenennungen und keine Datenänderungen erkennt.
\item Migration einspielen und prüfen, dass Modell und Datenbank übereinstimmen.
\item Im Review gleicht Julian Ammann Modell und Migration mit dem ER-Diagramm ab.
\end{enumerate}


In Produktion spielt ein eigener Container die Migrationen ein, bevor die API startet.

\subsubsection{Scheduler, E-Mail und Bilder}

Diese Teile entstehen in Sprint 2.

Der Scheduler (PSP 4.5, Kevin Jordan Taghu) arbeitet in zwei Schritten. Zuerst legt er fällige Benachrichtigungen als Zeilen an, danach versendet er sie. Ein eindeutiger Schlüssel aus Konto, Bezug, Termin, Typ und Intervall verhindert doppelte Zeilen. Beim Versand sperrt jede Instanz nur die Zeilen, die sie gerade bearbeitet. Auch wenn versehentlich zwei Scheduler laufen, geht keine Nachricht doppelt hinaus (QZ-05). Bleibt eine Nachricht länger als 15 Minuten im Versand hängen, protokolliert der Scheduler eine Warnung, statt sie erneut zu senden.

Der E-Mail-Versand (PSP 4.5, Kevin Jordan Taghu) liegt hinter einer Schnittstelle mit drei Umsetzungen: SMTP für Produktion, Mailpit als lokales Test-Postfach für Entwicklung und End-to-End-Tests und ein Aufzeichner für Unit-Tests. Tests und Demos schreiben nie echte Adressen an. Für Produktion stehen Brevo, Resend und Mailgun zur Wahl, als Rückfallebene der SMTP-Zugang eines Teammitglieds. Die Wahl fällt zu Beginn von Sprint 2.

Für Bilder (PSP 4.4, Anton Hirsch) sind JPEG, PNG und WebP bis 5 MB erlaubt. Das Format prüft das Backend am Dateiinhalt, nicht an der Endung. Jedes Bild wird neu kodiert, dabei fallen EXIF-Daten wie Aufnahmeort weg. Dateien liegen je Konto getrennt unter einem zufälligen Namen und sind nur mit gültiger Sitzung oder gültigem Teilen-Link abrufbar.

\subsection{Konfiguration des Frontends}

Das Frontend setzen Julian Ammann und Yin Yin Wu-Hanke gemeinsam um und prüfen sich gegenseitig.

\subsubsection{Datenfluss}

{\footnotesize
\begin{verbatim}
Browser --> SvelteKit (Node, :3000) --> FastAPI (:8000)
              |
              +- Seitenaufbau und Formulare: serverseitiger Aufruf des Backends
              |
              +- Aufrufe des Browsers auf /api/*: Weiterleitung an das Backend
\end{verbatim}
}

\subsubsection{Verbindliche Regeln}

\begin{enumerate}
\item Seiten laden ihre Daten serverseitig. Der API-Client wird pro Anfrage erzeugt und nicht global geteilt, weil sich sonst beim serverseitigen Rendern die Sitzungen verschiedener Nutzer vermischen könnten.
\item Schreibende Aktionen sind HTML-Formulare, die SvelteKit serverseitig verarbeitet. JavaScript verbessert sie nur, nötig ist es nicht, und alles lässt sich per Tastatur bedienen (QZ-04).
\item Fehlermeldungen des Backends zeigt das Frontend in verständlichem Deutsch am betroffenen Formular. Screenreader kündigen sie an. Eine zentrale Fehlerseite fängt unerwartete Fehler ab.
\item Typen stammen ausschließlich aus der generierten Typdatei. Sie wird nie von Hand bearbeitet.
\item Statuswerte der API übersetzt das Frontend ins Deutsche.
\end{enumerate}


\subsubsection{Routen und Sitzung}

\begin{tabularx}{\linewidth}{@{}>{\hsize=0.68\hsize\RaggedRight\arraybackslash}X >{\hsize=1.14\hsize\RaggedRight\arraybackslash}X >{\hsize=1.18\hsize\RaggedRight\arraybackslash}X@{}}
\toprule
\textbf{Bereich} & \textbf{Zweck} & \textbf{Schutz} \\
\midrule
\endhead
Anmeldung, Registrierung & Konto anlegen und anmelden & öffentlich. Angemeldete Nutzer werden zur Übersicht weitergeleitet \\
Übersicht, fachliche Seiten & alle kontogebundenen Funktionen des Geschenke-Managers & nur mit Sitzung, sonst Weiterleitung zur Anmeldung mit Rücksprungziel \\
Geteilte Ansicht & ausdrücklich freigegebene Personen und Geschenkideen anzeigen & öffentlich nur mit gültigem, nicht erratbarem Freigabe-Link; Nur-Lese-Zugriff \\
Konto & Konto anzeigen und löschen (F-17) & nur mit Sitzung \\
Abmeldung & Sitzung beenden, Cookie löschen & nur als Formularaktion \\
\texttt{/\allowbreak{}api/\allowbreak{}*} & Weiterleitung von Browser-Aufrufen an das Backend & Cookie wird durchgereicht \\
\bottomrule
\end{tabularx}

Bei jeder Anfrage prüft das Frontend serverseitig das Session-Cookie beim Backend und stellt den angemeldeten Nutzer allen Seiten bereit. Nach der Anmeldung leitet es nur auf eigene Pfade weiter, damit kein Link auf fremde Seiten umlenken kann (Open Redirect). Außerdem setzt es Sicherheits-Header gegen Einbettung in fremde Seiten und gegen falsch erkannte Dateitypen.

Die fachlichen Seiten entstehen in derselben Reihenfolge wie die Backend-Domänen. Das Frontend beginnt mit einer Domäne, sobald deren Endpunkte in der OpenAPI-Spezifikation stehen.

\subsection{Entwicklungsumgebung}

Jedes Teammitglied nutzt die eigene IDE. Damit alle mit denselben Werkzeugen arbeiten, pinnt das Werkzeug mise alle Versionen zentral im Repository:

\begin{tabularx}{\linewidth}{@{}l l@{}}
\toprule
\textbf{Werkzeug} & \textbf{Version} \\
\midrule
\endhead
Python & 3.14 \\
Node.js & 26.9.0 \\
pnpm & 12.5.1 \\
uv & 0.12.9 \\
Lefthook & 2.1.14 \\
gitleaks & 8.30.1 \\
\bottomrule
\end{tabularx}

Zusätzlich braucht jedes Teammitglied Docker mit Compose. Die Einrichtung eines neuen Rechners umfasst vier Schritte: Repository klonen, Werkzeuge über mise installieren, Abhängigkeiten und Git-Hooks mit einem gemeinsamen Setup-Befehl einrichten, lokale Datenbank als Container starten. Danach laufen Backend und Frontend mit automatischem Neuladen lokal.

Für wiederkehrende Aufgaben gibt es gemeinsame Befehle:

\begin{tabularx}{\linewidth}{@{}>{\hsize=0.76\hsize\RaggedRight\arraybackslash}X >{\hsize=1.24\hsize\RaggedRight\arraybackslash}X@{}}
\toprule
\textbf{Befehl} & \textbf{Zweck} \\
\midrule
\endhead
\texttt{mise run check} & Lint, Typprüfung und Tests für Backend und Frontend \\
\texttt{mise run openapi} & OpenAPI-Spezifikation exportieren und Frontend-Typen erzeugen \\
\texttt{mise run docs} & Dokumentation aus dem Code erzeugen \\
\texttt{mise run docker} & Container-Images bauen \\
\bottomrule
\end{tabularx}

Eine lokale Compose-Ergänzung öffnet die Ports für Datenbank, Backend und Frontend und startet im Entwicklungsmodus. Die Produktion verwendet sie nicht.

\subsection{Versionsverwaltung, Branching und Review}

\begin{itemize}
\item Der Branch \texttt{main} ist geschützt. Änderungen kommen nur per Pull Request mit mindestens einer Freigabe hinein. Force-Pushes und das Löschen des Branches sind gesperrt.
\item Für Container-, Workflow- und Deploy-Dateien ist Julian Ammann als Code Owner und Pflicht-Reviewer eingetragen, weil diese Dateien bestimmen, was auf dem Produktionsserver läuft.
\item Branches sind kurzlebig und heißen \texttt{feature/\allowbreak{}<domäne>-<thema>}, \texttt{fix/\allowbreak{}<thema>}, \texttt{docs/\allowbreak{}<thema>} oder \texttt{chore/\allowbreak{}<thema>}.
\item Commit-Nachrichten folgen Conventional Commits mit der Komponente als Scope, zum Beispiel \texttt{feat(backend): \ldots{}}, \texttt{fix(frontend): \ldots{}}, \texttt{docs(ms3): \ldots{}}.
\item Ein Pull Request umfasst eine Domäne oder ein Thema. Er wird zusammengeführt, wenn die CI grün ist, die Reviews laut Rollentabelle vorliegen und alle Anmerkungen bearbeitet sind.
\end{itemize}


\subsection{Lokale Qualitätssicherung über Git-Hooks}

\begin{tabularx}{\linewidth}{@{}l >{\hsize=1.34\hsize\RaggedRight\arraybackslash}X >{\hsize=0.66\hsize\RaggedRight\arraybackslash}X@{}}
\toprule
\textbf{Zeitpunkt} & \textbf{Prüfung} & \textbf{Komponente} \\
\midrule
\endhead
vor Commit & Ruff Lint mit Auto-Fix und Ruff Format & Backend \\
vor Commit & ESLint mit Auto-Fix und Prettier & Frontend \\
vor Commit & Neuerzeugung der Frontend-Typen, wenn sich Endpunkte oder Schemas ändern & Backend $\rightarrow$ Frontend \\
vor Commit & Secret-Scan der geänderten Dateien mit gitleaks & gesamt \\
vor Push & Typprüfung mit Pyright im Strict-Modus & Backend \\
vor Push & Prüfung auf fehlende Migrationen & Backend \\
vor Push & Typprüfung mit \texttt{svelte-check} & Frontend \\
\bottomrule
\end{tabularx}

Hooks lassen sich lokal umgehen. Die CI wiederholt deshalb alle Prüfungen.

\subsection{Continuous Integration}

GitHub Actions führt die Prüfungen nur für die geänderte Komponente aus. Bei reinen Dokumentationsänderungen meldet ein Platzhalter-Workflow die Pflichtprüfungen als erfolgreich, damit der Branchschutz nicht blockiert.

\begin{tabularx}{\linewidth}{@{}l >{\hsize=0.73\hsize\RaggedRight\arraybackslash}X >{\hsize=1.27\hsize\RaggedRight\arraybackslash}X@{}}
\toprule
\textbf{Workflow} & \textbf{Auslöser} & \textbf{Prüfungen} \\
\midrule
\endhead
Backend & Änderungen im Backend & Lint, Typprüfung, Migrationen gegen PostgreSQL 17 einspielen und auf Vollständigkeit prüfen, Tests und Abhängigkeitsscan \\
Frontend & Änderungen im Frontend & generierte Typen aktuell, Formatierung, Lint, Typprüfung, Unit- und Komponententests im Browser, Build und Abhängigkeitsscan \\
Docker & Änderungen an Backend oder Frontend & Build beider Images, Veröffentlichung in der GitHub Container Registry bei \texttt{main} und Versions-Tags \\
Secrets & jeder Push auf \texttt{main} und jeder Pull Request & Secret-Scan über die gesamte Historie \\
\bottomrule
\end{tabularx}

Die Frontend-CI prüft, dass die generierten Typen zur OpenAPI-Spezifikation passen. Eine automatische Prüfung der OpenAPI-Spezifikation gegen den Backend-Code ist im bestehenden Backend-Workflow noch nicht eingerichtet und wird vor MS 4 ergänzt. Der API-Vertrag ist damit erst nach dieser Ergänzung auf beiden Seiten abgesichert. Alle Workflows laufen mit Leserechten, nur der Docker-Workflow darf zusätzlich Images veröffentlichen. Die CI nutzt dieselben gepinnten Werkzeugversionen wie die lokale Umgebung.

\subsection{Abhängigkeiten}

\begin{itemize}
\item Alle Abhängigkeiten sind in Lockfiles fixiert. Die CI installiert nur aus den Lockfiles.
\item Entwicklungswerkzeuge sind von den Laufzeitabhängigkeiten getrennt und gelangen nicht in die Produktions-Images.
\item Dependabot prüft künftig wöchentlich Python- und npm-Pakete, GitHub Actions und Docker-Basis-Images und öffnet gebündelte Pull Requests. Die Konfiguration wird vor MS 4 auf \texttt{main} eingerichtet.
\item \texttt{pip-audit} und \texttt{pnpm audit} laufen in den bestehenden Workflows derzeit mit \texttt{continue-on-error}; Funde blockieren einen Merge daher noch nicht. Der Backend-Scan prüft zudem bisher die eigene Werkzeugumgebung von \texttt{pip-audit} statt der gesperrten Projektabhängigkeiten. Vor MS 4 werden beide Scans korrigiert und die Fehlerfortsetzungen entfernt. Im Frontend blockieren dann Befunde der Schweregrade hoch und kritisch einen Merge (\texttt{pnpm audit --audit-level=high}); mittlere und niedrige Befunde werden in der Bugliste mit Priorität und Zieltermin erfasst. \texttt{pip-audit} weist keine Schweregrade aus, deshalb blockiert im Backend jeder Befund. Begründete Ausnahmen dokumentiert die verantwortliche Person in einem Issue mit Risiko und Behebungsziel und hinterlegt sie per \texttt{--ignore-vuln}; Yin Yin Wu-Hanke prüft und genehmigt sie (QZ-06). Der Docker-Workflow baut derzeit Images, scannt sie aber noch nicht; ein Image-Scan wird vor MS 4 ergänzt.
\end{itemize}


\subsection{Konfiguration und Geheimnisse}

Die Anwendung liest ihre Konfiguration ausschließlich aus Umgebungsvariablen. Jede Variable ist im Code typisiert und beschrieben. Daraus entsteht automatisch die Konfigurationstabelle der Betriebsdokumentation.

\begin{tabularx}{\linewidth}{@{}>{\hsize=0.76\hsize\RaggedRight\arraybackslash}X >{\hsize=0.79\hsize\RaggedRight\arraybackslash}X >{\hsize=1.45\hsize\RaggedRight\arraybackslash}X@{}}
\toprule
\textbf{Variable} & \textbf{Standard} & \textbf{Bedeutung} \\
\midrule
\endhead
\texttt{APP\_\allowbreak{}ENV} & \texttt{development} & \texttt{development} aktiviert die API-Dokumentation und lesbare Logs, \texttt{production} JSON-Logs \\
\texttt{LOG\_\allowbreak{}LEVEL} & \texttt{INFO} & \texttt{DEBUG}, \texttt{INFO}, \texttt{WARNING}, \texttt{ERROR} \\
\texttt{DATABASE\_\allowbreak{}URL} & lokale Entwicklungs-DB & Verbindung zu PostgreSQL \\
\texttt{DEFAULT\_\allowbreak{}TIMEZONE} & \texttt{Europe/\allowbreak{}Berlin} & Zeitzone für Kalenderauswertungen \\
\texttt{UPLOADS\_\allowbreak{}DIR} & \texttt{/\allowbreak{}data/\allowbreak{}uploads} & Ablage der Bilddateien \\
\texttt{SESSION\_\allowbreak{}TTL\_\allowbreak{}DAYS} & \texttt{14} & Laufzeit einer Sitzung in Tagen \\
\texttt{SMTP\_\allowbreak{}*} & keiner & Zugang zum E-Mail-Anbieter \\
\texttt{API\_\allowbreak{}URL} & \texttt{http:/\allowbreak{}/\allowbreak{}localhost:8000} & Adresse des Backends aus Sicht des Frontends \\
\texttt{ORIGIN} & keiner & öffentliche Adresse des Frontends, in Produktion Pflicht \\
\bottomrule
\end{tabularx}

Regeln für Geheimnisse:

\begin{itemize}
\item \texttt{.env}-Dateien sind von der Versionierung ausgeschlossen. Die Vorlage \texttt{.env.example} enthält nur unkritische Entwicklungswerte.
\item gitleaks prüft vor jedem Commit und in der CI über die gesamte Historie.
\item Produktionsgeheimnisse wie Datenbank- und SMTP-Zugang liegen nicht im Projekt-Repository, sondern verschlüsselt im separaten Infrastruktur-Repository. Sie werden erst beim Deploy eingespielt.
\item Testkonten und Zugangsdaten für den Tutor übergibt das Team ausschließlich über Redmine.
\end{itemize}


\subsection{Teststrategie}

{\small
\begin{tabularx}{\linewidth}{@{}>{\hsize=0.58\hsize\RaggedRight\arraybackslash}X >{\hsize=0.89\hsize\RaggedRight\arraybackslash}X >{\hsize=1.79\hsize\RaggedRight\arraybackslash}X >{\hsize=0.74\hsize\RaggedRight\arraybackslash}X@{}}
\toprule
\textbf{Ebene} & \textbf{Werkzeug} & \textbf{Umfang} & \textbf{Verantwortlich} \\
\midrule
\endhead
Unit Backend & pytest mit Test-Uhr und Test-Postfach & Services, Planung der Benachrichtigungen, Vorschlagslogik & Entwickler der Domäne \\
Integration Backend & pytest, httpx, testcontainers & Kontogebundene Endpunkte mit Normalfall, Validierung und Fremdzugriff $\rightarrow$ 404. Freigabe-Endpunkte zusätzlich mit gültigem, ungültigem, abgelaufenem und widerrufenem Link; nur Lesen. Echte PostgreSQL-Datenbank, jeder Test wird danach zurückgerollt & Entwickler der Domäne \\
Berechtigung & wie Integration & Mandantentrennung je Ressource, Teilen-Link zeigt nur freigegebene Personen & Yin Yin Wu-Hanke \\
Unit Frontend & Vitest im Browser & Hilfsfunktionen und Komponenten & Julian Ammann \\
End-to-End & Playwright gegen den kompletten Stack & Kernabläufe: registrieren, Person anlegen, Idee $\rightarrow$ Beschenkung $\rightarrow$ verschenkt, Teilen-Link ohne Anmeldung öffnen, unzulässige Änderungen über den Link abweisen und Konto löschen. Chromium, Firefox und WebKit; Safari wird zusätzlich manuell geprüft (QZ-08) & Julian Ammann, Anton Hirsch \\
Last & Skript mit synthetischem Testbestand & Nachweis QZ-02, Messwerte gehen in den Testabschlussbericht & Anton Hirsch \\
\bottomrule
\end{tabularx}
}

\begin{itemize}
\item Ziel aus QZ-07 sind mindestens 70 \% Zeilenabdeckung im Backend. Der bestehende Backend-Workflow misst diese Abdeckung noch nicht; \texttt{pytest-cov} und der Grenzwert werden vor MS 4 in der CI eingerichtet.
\item Jeder Test nennt die Anforderungen, die er prüft. Sobald die Testkennungen in den Tests und im CI-Export verfügbar sind, entsteht aus jedem vollständigen Testlauf die Rückverfolgbarkeitsmatrix Anforderung $\leftrightarrow$ Test $\leftrightarrow$ Ergebnis als Nachweis für QZ-01.
\item Warnungen gelten in den Backend-Tests als Fehler.
\item End-to-End-Tests laufen nur auf \texttt{main} und vor Releases, weil sie Pull Requests sonst zu stark verlangsamen. Playwright prüft dabei WebKit als Annäherung an Safari, nicht den markengebundenen Safari-Browser; deshalb erfolgt die zusätzliche Safari-Prüfung manuell (siehe \href{https://playwright.dev/docs/browsers}{Playwright-Dokumentation zu Browsern}).
\end{itemize}


\subsection{Bereitstellung und Betrieb}

\subsubsection{Container}

\begin{tabularx}{\linewidth}{@{}l >{\hsize=0.65\hsize\RaggedRight\arraybackslash}X >{\hsize=1.35\hsize\RaggedRight\arraybackslash}X@{}}
\toprule
\textbf{Image} & \textbf{Basis} & \textbf{Merkmale} \\
\midrule
\endhead
Backend & \texttt{python:3.14-slim} & mehrstufiger Build, kein Root-Benutzer, nur Laufzeitabhängigkeiten. Dient auch für Migration und Scheduler \\
Frontend & \texttt{node:26-alpine} & mehrstufiger Build, kein Root-Benutzer \\
\bottomrule
\end{tabularx}

Die Images werden in der GitHub Container Registry mit Commit-Kennung und bei Releases mit der Versionsnummer veröffentlicht.

\subsubsection{Compose-Konfiguration}

Die Container-Konfiguration besteht aus drei Dateien:

\begin{tabularx}{\linewidth}{@{}l >{\hsize=1.00\hsize\RaggedRight\arraybackslash}X@{}}
\toprule
\textbf{Datei} & \textbf{Inhalt} \\
\midrule
\endhead
Basis & Dienste Datenbank, Migration, API, Scheduler und Frontend. Health-Checks, internes Netz, Volumes für Datenbank und Bilder. Keine offenen Ports \\
Lokal & Build aus dem Quellcode, offene Ports, Entwicklungsmodus \\
Produktion & Anbindung des Frontends an den Reverse Proxy mit der öffentlichen Domain, Speicher- und CPU-Limits für alle Dienste \\
\bottomrule
\end{tabularx}

Health-Checks steuern die Startreihenfolge. Die API startet erst nach erfolgreicher Migration, das Frontend erst, wenn die API bereit ist. Zwei Health-Endpunkte melden, ob der Prozess läuft und ob die Datenbank erreichbar ist.

\subsubsection{Hosting}

Das System läuft auf einem bestehenden, gehärteten Server eines Teammitglieds. Dort übernimmt Traefik Routing, TLS-Zertifikate über Let's Encrypt, HSTS, Sicherheits-Header und ein Rate-Limit am Eingang. Firewall, Monitoring und Log-Sammlung stellt ebenfalls der Server. Das Projekt liefert nur seine Container-Konfiguration.

Julian Ammann stößt das Deploy im separaten Infrastruktur-Repository an, sobald die Images gebaut sind. Deployt wird immer ein fester Commit mit dem dazu gebauten Image, nie der jeweils aktuelle Stand von \texttt{main}. Weil das Deploy dort angestoßen wird, liegen im öffentlichen Projekt-Repository keine Deploy-Zugangsdaten. Vor jedem Deploy lehnen automatische Prüfungen offene Ports, Einbindungen von Host-Verzeichnissen und fremde Netze ab.

Schon in Sprint 0 ist die leere Anwendung mit Health-Endpunkt unter HTTPS erreichbar, damit Probleme bei der Bereitstellung nicht erst kurz vor MS 4 auffallen.

Das System verarbeitet nur synthetische Testdaten, ein externes Backup ist deshalb nicht vorgesehen. Nach einem Datenverlust spielt ein Seed-Skript Testkonten und Testbestand neu ein. Das Skript entsteht vor der Bereitstellung für den Tutor.

\subsection{Definition of Done}

Eine Funktion ist fertig, wenn alle Punkte erfüllt sind:

\paragraph{Backend}

\begin{itemize}
\item[$\square$] Tabelle entspricht dem verbindlichen Datenmodell, Migration erzeugt, gegengelesen und vollständig
\item[$\square$] jede Tabelle und Spalte hat einen Datenbankkommentar
\item[$\square$] Service filtert jede Abfrage nach dem Besitzer, fachliche Fehler nutzen die gemeinsame Fehlerklasse
\item[$\square$] Kontogebundener Endpunkt erfordert Anmeldung; öffentliche Freigabe-Endpunkte erlauben ausschließlich den über einen gültigen Link freigegebenen Nur-Lese-Zugriff. Jeder Endpunkt hat Operation-ID, Zusammenfassung und Beschreibung. Schema-Felder haben Beschreibung und Beispiel
\item[$\square$] Listen sind paginiert
\item[$\square$] Tests für Normalfall, Validierung und Fremdzugriff $\rightarrow$ 404, jeweils mit Anforderungs-ID
\item[$\square$] OpenAPI-Spezifikation und generierte Dokumentation sind aktualisiert und committet
\end{itemize}


\paragraph{Frontend}

\begin{itemize}
\item[$\square$] Daten werden serverseitig geladen, Änderungen laufen über Formularaktionen
\item[$\square$] Fehler sind sichtbar und werden von Screenreadern angekündigt. Die Seite ist ohne JavaScript bedienbar
\item[$\square$] Typen stammen nur aus der generierten Typdatei. Lint, Typprüfung und Tests sind grün
\end{itemize}


\paragraph{Gemeinsam}

\begin{itemize}
\item[$\square$] CI grün, Review laut Rollentabelle freigegeben
\item[$\square$] betroffene Anforderung steht in der Rückverfolgbarkeitsmatrix als bestanden
\end{itemize}


\section{Konfiguration der Liefergegenstände}

\subsection{Liefergegenstände und Zuständigkeiten}

Jeder Liefergegenstand erhält eine Kennung, eine verantwortliche und eine prüfende Person sowie einen festen Ablageort. Die prüfende Person ist nie die verantwortliche.

{\small
\begin{tabularx}{\linewidth}{@{}l >{\hsize=1.45\hsize\RaggedRight\arraybackslash}X l >{\hsize=0.93\hsize\RaggedRight\arraybackslash}X >{\hsize=0.87\hsize\RaggedRight\arraybackslash}X >{\hsize=0.99\hsize\RaggedRight\arraybackslash}X >{\hsize=0.75\hsize\RaggedRight\arraybackslash}X@{}}
\toprule
\textbf{ID} & \textbf{Liefergegenstand} & \textbf{MS} & \textbf{Format} & \textbf{Ablage} & \textbf{Verantwortlich} & \textbf{Prüfend} \\
\midrule
\endhead
LG-01 & Konfiguration der Softwareentwicklung und Qualitätsplanung & 3 & Markdown $\rightarrow$ PDF & \texttt{docs/\allowbreak{}} & Julian Ammann & Kevin Jordan Taghu \\
LG-02 & Programmcode Backend & 4 & Git-Repository & \texttt{backend/\allowbreak{}} & Kevin Jordan Taghu, Anton Hirsch (Anmeldung: Julian Ammann) & Julian Ammann \\
LG-03 & Programmcode Frontend & 4 & Git-Repository & \texttt{frontend/\allowbreak{}} & Julian Ammann, Yin Yin Wu-Hanke & gegenseitiges Review \\
LG-04 & Lauffähiges System (Link) & 4 & Container-Images, URL & GitHub Container Registry, Produktionsserver & Julian Ammann & Anton Hirsch \\
LG-05 & Benutzerhandbuch & 4 & Markdown $\rightarrow$ PDF mit Screenshots & \texttt{docs/\allowbreak{}ms4/\allowbreak{}} & Yin Yin Wu-Hanke & Anton Hirsch \\
LG-06 & Fachliche Dokumentation: Prozesse, Konzepte, Geschäftsregeln & 4 & Markdown $\rightarrow$ PDF & \texttt{docs/\allowbreak{}ms4/\allowbreak{}} & Anton Hirsch & Kevin Jordan Taghu \\
LG-07 & Technische Dokumentation: Architektur, Komponenten, Schnittstellen, Datenbank & 4 & Markdown, teilweise generiert & \texttt{docs/\allowbreak{}ms4/\allowbreak{}}, \texttt{docs/\allowbreak{}generated/\allowbreak{}} & Julian Ammann (Architektur, Frontend), Kevin Jordan Taghu (Backend) & Anton Hirsch \\
LG-08 & Betriebsdokumentation: Installation, Konfiguration, Admin-Account & 4 & Markdown, teilweise generiert & \texttt{deploy/\allowbreak{}}, \texttt{docs/\allowbreak{}generated/\allowbreak{}} & Julian Ammann & Kevin Jordan Taghu \\
LG-09 & Testabschlussbericht mit Testfällen und Testprotokollen & 4 & Markdown $\rightarrow$ PDF, CI-Protokolle & \texttt{docs/\allowbreak{}ms4/\allowbreak{}}, \texttt{docs/\allowbreak{}generated/\allowbreak{}} & Anton Hirsch & Yin Yin Wu-Hanke \\
LG-10 & Liste der Testkonten und Zugangsdaten & 4 & PDF & nur Redmine, nie im Repository & Yin Yin Wu-Hanke & Julian Ammann \\
LG-11 & Ergebnispräsentation: Projektablauf mit Teilergebnissen, Liefergegenständen und Reflexion ihrer Erstellung; technischer Überblick, Testabschlussbericht, Demo und Lessons Learned & 5 & Video oder Link, höchstens 20 Minuten; Demo höchstens 7 Minuten & \texttt{docs/\allowbreak{}ms5/\allowbreak{}} & Kevin Jordan Taghu & Anton Hirsch \\
LG-12 & Gemeinsamer Projektbericht & 6 & PDF & \texttt{docs/\allowbreak{}ms6/\allowbreak{}} & Kevin Jordan Taghu & alle \\
\bottomrule
\end{tabularx}
}

\subsection{Gliederung, Inhalte und Umfang}

\begin{tabularx}{\linewidth}{@{}l >{\hsize=1.20\hsize\RaggedRight\arraybackslash}X >{\hsize=0.80\hsize\RaggedRight\arraybackslash}X@{}}
\toprule
\textbf{ID} & \textbf{Gliederung und Inhalte} & \textbf{Umfang} \\
\midrule
\endhead
LG-01 & Konfiguration der Softwareentwicklung, Konfiguration der Liefergegenstände, Qualitätsplanung (Gliederung dieses Dokuments) & alle in der Aufgabenstellung zu MS 3 genannten Inhalte \\
LG-02 & Backend nach der Repository-Struktur: Migrationen, Router, Kern, Modelle, Schemas, Services, Tests & alle MUSS-Anforderungen mit Tests; Definition of Done je Funktion erfüllt \\
LG-03 & Frontend nach der Repository-Struktur: API-Client mit generierten Typen, Seiten, Proxy für \texttt{/\allowbreak{}api/\allowbreak{}*}, Tests & alle Oberflächen zu den MUSS-Anforderungen \\
LG-04 & Öffentliche Adresse, Container-Images mit Versionsnummer, Health-Endpunkte & im Browser ohne lokale Installation prüfbar \\
LG-05 & Einstieg und Anmeldung, Personen und Anlässe, Geschenkideen und Beschenkungen, Aufgaben und Notizen, Benachrichtigungen, Teilen-Links, Konto löschen, häufige Fragen. Jeder Ablauf mit Schrittfolge und Screenshot & ein Kapitel je Kernablauf aus den End-to-End-Tests \\
LG-06 & Fachliche Prozesse mit Sequenzdiagrammen, fachliche Konzepte (Person, Anlass, Geschenk/Idee, Beschenkung, Teilen-Link), Geschäftsregeln mit Verweis auf die Anforderungs-IDs, Glossar & jede MUSS-Anforderung mindestens einer Regel oder einem Prozess zugeordnet \\
LG-07 & Architekturüberblick und Architekturentscheidungen, Komponenten von Backend, Frontend und Scheduler, Schnittstellen (OpenAPI-Dokumentation), Datenbank (ER-Diagramm, Tabellen- und Spaltenbeschreibungen) & Schnittstellen und Datenbank vollständig generiert, übrige Kapitel von Hand \\
LG-08 & Voraussetzungen, Installation aus leerem Zustand, Konfiguration (Umgebungsvariablen), Admin-Account, Aktualisierung, Neuaufbau des Testbestands über das Seed-Skript & reicht einer nicht beteiligten Person für die Installation (QZ-08) \\
LG-09 & Testgegenstand und Testumgebung, Teststrategie, Testfälle je Ebene, Rückverfolgbarkeitsmatrix, Testprotokolle, Messwerte zu QZ-01 bis QZ-09, offene Fehler aus der Bugliste, Bewertung & alle Qualitätsziele mit Messwert und Ergebnis \\
LG-10 & Je Konto Zweck, Anmeldename und Passwort; Admin-Account gesondert gekennzeichnet & alle für die Prüfung durch den Tutor benötigten Konten \\
LG-11 & Projektablauf mit Teilergebnissen, Liefergegenständen und Reflexion; technischer Überblick; Überblick Testabschlussbericht; Demo; Lessons Learned & höchstens 20 Minuten, Demo höchstens 7 Minuten \\
LG-12 & Titelblatt mit Angaben aller Mitglieder und Seitenbereichen, Einleitung, Projektverlauf MS 0 bis MS 5 mit Ergebnissen, Reflexion und Fazit & je Person ein zusammenhängender, namentlich zugeordneter Textteil von 7 bis 10 Seiten. Jedes Mitglied reicht den Bericht gesondert in Turnitin ein \\
\bottomrule
\end{tabularx}

\subsection{Dokumentation aus dem Code}

MS 4 verlangt acht Liefergegenstände, davon fünf Dokumente. Ein großer Teil davon entsteht direkt beim Entwickeln, sonst müsste das Team alles in der letzten Woche schreiben. Nach den Regeln zur Dokumentation im Code tragen Endpunkte, Tabellen, Konfiguration und Tests ihre Beschreibung von Anfang an. Ein gemeinsamer Befehl erzeugt daraus die Dokumente.

{\small
\begin{tabularx}{\linewidth}{@{}>{\hsize=1.10\hsize\RaggedRight\arraybackslash}X >{\hsize=1.24\hsize\RaggedRight\arraybackslash}X >{\hsize=1.02\hsize\RaggedRight\arraybackslash}X >{\hsize=0.64\hsize\RaggedRight\arraybackslash}X@{}}
\toprule
\textbf{Liefergegenstand} & \textbf{Quelle im Code} & \textbf{Entstehung} & \textbf{Anteil aus Code} \\
\midrule
\endhead
Technische Doku: Schnittstellen & Beschreibungen an Endpunkten und Schema-Feldern & generiert, dazu OpenAPI-Spezifikation und interaktive API-Doku & fast vollständig \\
Technische Doku: Datenbank & Kommentare an Tabellen und Spalten, ER-Diagramm & generiert & weitgehend \\
Technische Doku: Architektur, Komponenten & Architekturentscheidungen aus diesem Dokument & von Hand & teilweise \\
Betriebsdoku: Konfiguration & Beschreibungen der Umgebungsvariablen & generiert & vollständig \\
Betriebsdoku: Installation, Admin-Account & Compose-Dateien, Seed-Skript & von Hand & teilweise \\
Testabschlussbericht & Anforderungs-IDs an Tests, Testläufe, Abdeckung & generiert aus dem Testlauf, ergänzt um Bewertung & weitgehend \\
Fachliche Doku: Geschäftsregeln & Beschreibungen der Services, Fehlermeldungen, Testnamen & von Hand & teilweise \\
Fachliche Doku: Prozesse, Konzepte & keine & von Hand, mit Sequenzdiagrammen und Glossar & kaum \\
Benutzerhandbuch & Screenshots aus den End-to-End-Tests & Text von Hand, Screenshots automatisch & kaum \\
\bottomrule
\end{tabularx}
}

Die generierten Dateien liegen unter \texttt{docs/\allowbreak{}generated/\allowbreak{}} und werden committet, damit man sie ohne Werkzeuge lesen und Änderungen im Pull Request sehen kann. Von Hand bearbeitet werden sie nie. Die Screenshots für das Benutzerhandbuch erzeugen die End-to-End-Tests bei jedem Lauf neu, sodass sie nach Änderungen an der Oberfläche aktuell bleiben.

\subsection{Identifikation und Versionierung}

\begin{itemize}
\item Dokumente liegen unter \texttt{docs/\allowbreak{}} und heißen \texttt{ms<Nummer>\_\allowbreak{}<thema>.md}. Ab MS 4 erhält jeder Meilenstein einen Unterordner \texttt{docs/\allowbreak{}ms<Nummer>/\allowbreak{}}. Dateinamen sind durchgehend klein geschrieben.
\item PDF-Abgabefassungen werden aus den Markdown- oder LaTeX-Quellen erzeugt und unter \texttt{docs/\allowbreak{}pdf/\allowbreak{}} abgelegt. Änderungen erfolgen immer in der Quelle, nie im PDF.
\item Die Software folgt Semantic Versioning, aktuell \texttt{0.1.0}. Die Abgabe zu MS 4 erhält die Version \texttt{1.0.0}. Das Setzen der Version löst Build und Veröffentlichung der Images mit derselben Versionsnummer aus. Korrekturen nach Tutorfeedback erhöhen die letzte Stelle, etwa \texttt{1.0.1}.
\item Jeder Abgabestand erhält einen Git-Tag \texttt{ms<Nummer>-abgabe}. Über den Tag lässt sich der eingereichte Stand jederzeit wiederherstellen, auch wenn \texttt{main} weiterentwickelt wird.
\end{itemize}


\subsection{Baselines je Meilenstein}

\begin{tabularx}{\linewidth}{@{}l >{\hsize=1.22\hsize\RaggedRight\arraybackslash}X >{\hsize=0.78\hsize\RaggedRight\arraybackslash}X@{}}
\toprule
\textbf{Baseline} & \textbf{Inhalt} & \textbf{Kennzeichnung} \\
\midrule
\endhead
MS 3 & dieses Dokument, Datenmodell, Entwicklungsumgebung, CI, Grundgerüst mit Health-Endpunkten & Tag \texttt{ms3-abgabe} \\
MS 4 & Programmcode, Container-Images, Dokumentation LG-05 bis LG-09, generierte Dokumente & Tags \texttt{v1.0.0} und \texttt{ms4-abgabe}, Images \texttt{1.0.0} \\
MS 5 & Präsentation auf Basis der MS-4-Baseline & Tag \texttt{ms5-abgabe} \\
MS 6 & Projektbericht & Tag \texttt{ms6-abgabe} \\
\bottomrule
\end{tabularx}

Nach einer Baseline ändert sich der abgegebene Stand nur über einen neuen Pull Request und eine neue Version.

\subsection{Änderungs- und Freigabeprozess}

\begin{enumerate}
\item Die verantwortliche Person erstellt oder ändert den Liefergegenstand in einem eigenen Branch.
\item Sie öffnet einen Pull Request gegen \texttt{main}. Die prüfende Person prüft Inhalt, Form und Vollständigkeit gegen die Qualitätsziele und die Definition of Done.
\item Die CI muss grün sein.
\item Nach der Freigabe wird der Pull Request zusammengeführt, die Abgabefassung erzeugt und der Stand getaggt.
\item Die Projektleitung lädt die Abgabe in Redmine hoch, setzt das Meilensteinticket auf \textbf{Feedback} und weist es dem Tutor zu.
\item Das Tutorfeedback steht am Redmine-Ticket. Notwendige Korrekturen durchlaufen erneut die Schritte 1 bis 5. Der Tutor nimmt mit dem Status \textbf{closed} ab.
\end{enumerate}


Die letzten zwei Tage vor jeder Abgabe sind für Review, Korrektur und Bereitstellung reserviert.

\subsection{Prüfverfahren für Liefergegenstände}

\begin{tabularx}{\linewidth}{@{}>{\hsize=0.62\hsize\RaggedRight\arraybackslash}X >{\hsize=1.38\hsize\RaggedRight\arraybackslash}X@{}}
\toprule
\textbf{Liefergegenstand} & \textbf{Prüfung} \\
\midrule
\endhead
Programmcode & CI grün (Lint, Typen, Tests, Schwachstellen- und Secret-Scan), Review laut Rollentabelle, Definition of Done erfüllt \\
Lauffähiges System & Installation aus leerem Zustand, Health-Endpunkte, End-to-End-Abläufe gegen die Produktionsadresse, Anmeldung mit allen Testkonten \\
Generierte Dokumente & aus dem Release-Stand erzeugt, Änderungen im Pull Request geprüft \\
Handgeschriebene Dokumente & eine nicht beteiligte Person liest den Text gegen den Code und gibt ihn im Pull Request frei \\
Testabschlussbericht & Zahlen zu Testanzahl, Abdeckung und Antwortzeiten nur aus echten Läufen des Release-Stands, Matrix vollständig \\
Testkonten & Anmeldung mit jedem Konto vor der Übergabe, keine Zugangsdaten im Repository \\
\bottomrule
\end{tabularx}

\subsection{Schutz vertraulicher Liefergegenstände}

Das Repository ist öffentlich. Zugangsdaten, Testkonten und Produktionsgeheimnisse werden deshalb nie versioniert. Testkonten erhält der Tutor nur über Redmine (LG-10), Produktionsgeheimnisse liegen im separaten Infrastruktur-Repository. Der Prototyp verarbeitet ausschließlich synthetische Testdaten. gitleaks fängt lokal und in der CI Geheimnisse ab, die versehentlich in einen Commit geraten.

\section{Qualitätsplanung}

\subsection{Qualitätsziele}

Die Qualitätsziele konkretisieren die Qualitätsanforderungen Q-01 bis Q-09 aus MS 1 und sind nach den Merkmalen der ISO/IEC 25010 gegliedert. Die Kennungen QZ ergänzen die ursprünglichen Kennungen Q; sie ersetzen diese nicht. Jedes Ziel ist messbar und einem Prüfverfahren, einem Zeitpunkt und einer verantwortlichen Person zugeordnet.

\begin{tabularx}{\linewidth}{@{}>{\hsize=0.89\hsize\RaggedRight\arraybackslash}X >{\hsize=1.11\hsize\RaggedRight\arraybackslash}X@{}}
\toprule
\textbf{Anforderung aus MS 1} & \textbf{Konkretisierung in MS 3} \\
\midrule
\endhead
Q-01 Mandantentrennung & QZ-06 Sicherheit \\
Q-02 Benachrichtigungszuverlässigkeit & QZ-05 Zuverlässigkeit \\
Q-03 Sicherheit & QZ-06 Sicherheit \\
Q-04 Datenschutz & QZ-06 Sicherheit \\
Q-05 Leistung & QZ-02 Leistungseffizienz \\
Q-06 Wartbarkeit & QZ-07 Wartbarkeit \\
Q-07 Barrierefreiheit & QZ-04 Benutzbarkeit und Barrierefreiheit \\
Q-08 Zeitkorrektheit & QZ-09 Zeitkorrektheit \\
Q-09 Nachvollziehbarkeit & QZ-01 Rückverfolgbarkeit der Anforderungen sowie Dokumentations- und Freigaberegeln für Entscheidungen und Liefergegenstände \\
\bottomrule
\end{tabularx}

QZ-03 (API-Vertrag und HTML-Konformität) und QZ-08 (Übertragbarkeit und Browserdarstellung) ergänzen die MS-1-Ziele.

{\small
\begin{tabularx}{\linewidth}{@{}l >{\hsize=0.68\hsize\RaggedRight\arraybackslash}X >{\hsize=1.67\hsize\RaggedRight\arraybackslash}X >{\hsize=1.16\hsize\RaggedRight\arraybackslash}X >{\hsize=0.48\hsize\RaggedRight\arraybackslash}X@{}}
\toprule
\textbf{ID} & \textbf{Merkmal} & \textbf{Qualitätsziel} & \textbf{Prüfverfahren} & \textbf{Prüfer} \\
\midrule
\endhead
QZ-01 & Funktionale Eignung & Alle MUSS Anforderungen sind umgesetzt und durch mindestens einen Testfall abgedeckt; 100 \% der zugehörigen Testfälle sind bestanden. Berechnungen sind korrekt, z. B. enthält die Geburtstagsübersicht des Folgemonats genau die betroffenen Personen. Die Umwandlung einer Geschenkidee in ein Geschenk erfolgt in einem Schritt. & Rückverfolgbarkeitsmatrix Anforderung <--> Testfall; Unit-, Integrations- und Systemtests & Yin Yin Wu-Hanke \\
QZ-02 & Leistungseffizienz & Listen mit 100 Personen, 1.000 Geschenkideen und 1.000 Beschenkungen werden in der festgehaltenen lokalen Testumgebung innerhalb von 2 Sekunden angezeigt (Q-05). & Nach einem Aufwärmlauf 30 Messungen vom Aufruf bis zur Anzeige der Liste; der 95. Perzentilwert muss höchstens 2 Sekunden betragen. Testgerät und Laufbedingungen werden im Testprotokoll festgehalten. & Julian Ammann \\
QZ-03 & Kompatibilität & Der Frontend-Client entspricht der OpenAPI-Spezifikation des Backends (typisierter Client, keine Typfehler). Der HTML-Export ist valides HTML. & Typprüfung in CI, W3C-Validator & Anton Hirsch, Kevin Jordan Taghu \\
QZ-04 & Benutzbarkeit und Barrierefreiheit & Eine Geschenkidee kann ohne Anleitung in höchstens 3 Interaktionen und unter 30 Sekunden erfasst werden. Löschvorgänge erfordern eine Bestätigung, ungültige Eingaben erzeugen verständliche Fehlermeldungen, wesentliche Abläufe funktionieren per Tastatur und Zustände werden nicht nur durch Farbe vermittelt. Lighthouse-Accessibility-Wert $\geq$ 90. & Drei externe Testpersonen führen dieselbe Erfassungsaufgabe aus; das Ziel gilt als erreicht, wenn alle drei die Erfassung innerhalb der vorgegebenen Zeit und Interaktionen ohne Hilfe abschließen. Zeit und Interaktionen werden einzeln protokolliert. Zusätzlich Testfälle für ungültige Eingaben und Tastaturbedienung sowie Lighthouse-Audit. & Anton Hirsch, Yin Yin Wu-Hanke \\
QZ-05 & Zuverlässigkeit & Benachrichtigungen werden zum vorgesehenen Termin erzeugt. Wiederholte oder parallele Scheduler-Läufe erzeugen dabei keine zusätzlichen Versandaufträge. Unklare Versandergebnisse werden nicht automatisch erneut versendet, sondern protokolliert. Ein Ausfall des E-Mail-Versands beeinträchtigt die übrigen Funktionen nicht. Container starten nach einem Neustart automatisch. & Scheduler-Test mit simulierter Systemzeit, Integrationstest mit simuliertem Ausfall, Neustart-Test & Yin Yin Wu-Hanke \\
QZ-06 & Sicherheit & Nutzende greifen ausschließlich auf eigene Daten zu; die einzige Ausnahme sind ausdrücklich freigegebene Inhalte über einen gültigen Freigabe-Link. Passwörter und Tokens werden nur gehasht gespeichert. Teilen-Links sind nicht erratbar, optional befristet und jederzeit deaktivierbar. Keine Secrets im Repository. Es werden ausschließlich synthetische Daten verarbeitet, Vorschläge nutzen keine personenbezogenen Daten anderer Konten. & Autorisierungstests je Endpunkt, pnpm audit, pip-audit, Image-Scan (vor MS 4 einzurichten), Secret-Scan, Checkliste Datenschutz & Kevin Jordan Taghu, Anton Hirsch \\
QZ-07 & Wartbarkeit & Linter und Typprüfung laufen ohne Fehler. Jeder Merge auf main erfolgt nach Code-Review. Die Zeilenabdeckung der Backend-Unit-Tests $\geq$ 70 \% und jeder API-Endpunkt besitzt mindestens einen automatisierten Test. Frontend und Backend sind getrennt bau- und testbar. & ESLint, svelte-check, Ruff, pytest-cov, Branch Protection, Endpunktliste gegen Testliste & Julian Ammann \\
QZ-08 & Übertragbarkeit & Das Gesamtsystem startet aus leerem Zustand mit \texttt{docker compose up} anhand der Betriebsdokumentation. Die Kernfunktionen sind in aktuellen Versionen von Chrome, Firefox und Safari sowie auf Smartphone-Displaybreite nutzbar. & Deployment-Test auf sauberer Umgebung durch nicht beteiligte Person, Playwright mit Chromium, Firefox und WebKit sowie manueller Safari-Prüfung & Kevin Jordan Taghu, Yin Yin Wu-Hanke \\
QZ-09 & Zeitkorrektheit & Technische Zeitpunkte werden in UTC gespeichert; Geburtstage, Anlasstermine und Verschenkdaten bleiben reine Kalenderdaten. Benachrichtigungen berücksichtigen \texttt{Europe/\allowbreak{}Berlin}; Zeitumstellungen verschieben keinen Anlass auf einen anderen Kalendertag. & Unit- und Integrationstests mit fester Test-Uhr, Jahreswechseln, Schaltjahren und Zeitumstellungen & Kevin Jordan Taghu \\
\bottomrule
\end{tabularx}
}

\subsection{Qualität des Softwaresystems}

Die Qualitätsziele QZ-01 bis QZ-09 werden über die in der Konfiguration der Softwareentwicklung beschriebenen Maßnahmen erreicht: Teststrategie, lokale Git-Hooks, Continuous Integration, Code-Review und Definition of Done. Jede Änderung durchläuft dieselben Prüfschritte, bevor sie auf \texttt{main} gelangt.

\begin{tabularx}{\linewidth}{@{}>{\hsize=0.56\hsize\RaggedRight\arraybackslash}X >{\hsize=1.65\hsize\RaggedRight\arraybackslash}X >{\hsize=0.79\hsize\RaggedRight\arraybackslash}X@{}}
\toprule
\textbf{Zeitpunkt} & \textbf{Prüfung} & \textbf{Nachweis} \\
\midrule
\endhead
bei jedem Commit & Formatierung, Linting und Secret-Scan über die Git-Hooks & Ausgabe der Hooks \\
bei jedem Pull Request & Lint, Typprüfung, Unit- und Integrationstests, Build, Schwachstellenscan und mindestens eine Review-Freigabe; nach Aktivierung der Richtlinie blockieren Befunde nach den Regeln unter „Abhängigkeiten“ & grüne CI und Freigabe im Pull Request \\
je Sprint & funktionale Abnahme der umgesetzten Anforderungen (QZ-01) und, sobald in der CI eingerichtet, Messung der Testabdeckung (QZ-07) & Rückverfolgbarkeitsmatrix, Abdeckungsbericht \\
vor MS 4 & End-to-End-Abläufe über mehrere Browser (QZ-08), Last- und Antwortzeitmessung (QZ-02), Barrierefreiheit (QZ-04), Sicherheitsprüfung (QZ-06) und Zeitkorrektheit (QZ-09) & Testabschlussbericht \\
\bottomrule
\end{tabularx}

Neben den automatischen Prüfungen werden Artefakte statisch geprüft. Die Konkretisierung der Anforderungen prüft das Team vor der Umsetzung (siehe Anforderungsmanagement). Architekturentscheidungen werden vor der Umsetzung als GitHub Issue mit Begründung und betrachteten Alternativen festgehalten und von einer zweiten Person freigegeben; Änderungen an Container-, Workflow- und Deploy-Konfiguration prüft Julian Ammann. Testfälle werden im selben Pull Request wie der Code reviewt. Die prüfende Person kontrolliert dabei, ob sie die Akzeptanzkriterien der genannten Anforderung abdecken und neben dem Normalfall auch Validierung und Fremdzugriff prüfen.

Zentrale Nachweise sind die Rückverfolgbarkeitsmatrix Anforderung $\leftrightarrow$ Test $\leftrightarrow$ Ergebnis (QZ-01), die in der CI gemessene Zeilenabdeckung von mindestens 70 \% im Backend (QZ-07) sowie die grüne Continuous Integration mit bestandenem Review als Voraussetzung für jede Zusammenführung. Die abschließenden Messwerte stammen ausschließlich aus echten Läufen des Release-Stands und fließen in den Testabschlussbericht.

Wird ein Qualitätsziel verfehlt, wird die betroffene Änderung nicht zusammengeführt. Die Ursache wird behoben oder der Umfang wird angepasst, bevor weitergearbeitet wird.

\subsection{Qualität der Liefergegenstände}

Neben dem Softwaresystem werden auch die übrigen Liefergegenstände geplant geprüft, darunter Dokumente, Videos, Programmcode, das lauffähige System, der Testabschlussbericht und die Testkonten. Grundlage sind die Rollentabelle, die Prüfverfahren für Liefergegenstände und die Definition of Done aus der Konfiguration.

Für jeden Liefergegenstand gilt:

\begin{itemize}
\item Es sind eine verantwortliche und eine prüfende Person festgelegt. Die erstellende Person gibt ihren eigenen Liefergegenstand nicht allein frei.
\item Die Prüfung erfolgt vor der Abgabe anhand des festgelegten Prüfverfahrens. Das Ergebnis wird im Pull Request oder am zugehörigen Redmine-Ticket festgehalten.
\item Jeder Meilenstein hat einen eindeutig gekennzeichneten Abgabestand, auf den sich die Prüfung bezieht.
\end{itemize}


Für Dokumente und Präsentationen gelten zusätzliche Prüfpunkte:

\begin{tabularx}{\linewidth}{@{}>{\hsize=0.63\hsize\RaggedRight\arraybackslash}X >{\hsize=1.37\hsize\RaggedRight\arraybackslash}X@{}}
\toprule
\textbf{Merkmal} & \textbf{Prüfung} \\
\midrule
\endhead
Vollständigkeit & Alle für den Meilenstein geforderten Inhalte sind enthalten. \\
Fachliche Richtigkeit & Eine nicht beteiligte Person liest den Text gegen den aktuellen Stand von Code und Datenmodell. \\
Form & Einheitliche Struktur, Rechtschreibung, Format und Titelblatt werden geprüft. \\
Nachvollziehbarkeit & Aussagen sind mit Anforderungen, Artefakten oder dem Repository verknüpft. \\
\bottomrule
\end{tabularx}

Der Projektbericht (MS 6) wird zusätzlich über Turnitin eingereicht. Jedes Mitglied verantwortet seinen namentlich zugeordneten Textteil. Generierte Dokumente werden immer aus dem geprüften Release-Stand erzeugt, damit Dokumentation und System zusammenpassen.

\end{document}
